\documentclass[11pt]{article}

\usepackage[a4paper,margin=1in]{geometry}
\usepackage{amsmath}
\usepackage{amssymb}
\usepackage{booktabs}
\usepackage{enumitem}
\usepackage[colorlinks=true,linkcolor=blue,urlcolor=blue]{hyperref}
\usepackage{parskip}

\title{Deep-Learning Segmentation of Fibers in SEM Micrographs: Methods}
\author{}
\date{}

\begin{document}

\maketitle

\begin{abstract}
We trained a convolutional encoder--decoder network to perform binary semantic
segmentation of fiber structures in grayscale scanning electron microscopy
(SEM) images of sampled air-filter deposits. The model is intended to assist
-- not silently replace -- the manual fiber-counting workflow used in
occupational asbestos-exposure assessment (filter loading, WHO-criteria fiber
counts, and airborne concentration / time-weighted-average exposure
reporting). Training and evaluation are fully config-driven and tracked in
MLflow, so every number reported below traces back to an exact configuration
and code state.
\end{abstract}

\section{Dataset and splits}

The training corpus consists of 298 image/mask pairs of SEM micrographs (one
image excluded for a missing corresponding mask). Each image is a
single-channel grayscale micrograph; each mask is a binary image where a
pixel value greater than zero denotes fiber and zero denotes background. Mask
files are matched to images by filename pattern (\texttt{\{stem\}\_mask.tif}).

Pairs are split \textbf{by source image}, not by tile, into 70\% train / 15\%
validation / 15\% test (209 / 45 / 44 images) using a seeded random shuffle
(seed 42). Splitting by image before tiling is deliberate: it prevents tiles
from the same source micrograph from leaking across train/validation/test,
which would otherwise inflate validation/test performance.

\section{Preprocessing and tiling}

\paragraph{Contrast normalization.}
Each source image is normalized once, at full resolution, by clipping to its
1st and 99.5th percentile intensity and rescaling to $[0, 1]$. Percentile
(rather than min--max) normalization is used because SEM images commonly
contain a small number of very bright outlier pixels (charging artifacts,
detector saturation) that would otherwise dominate a naive min--max stretch.
This normalized array is cached to disk (memory-mapped \texttt{.npy}) so
every tile drawn from a given image shares an identical contrast stretch, and
the same normalization is applied verbatim at inference time.

\paragraph{Tiling.}
Each normalized image is decomposed into overlapping square tiles via a
sliding window (default $512\times512$~px, 384~px stride -- 25\% overlap).
The final row/column of tiles is snapped to the image edge so the entire
image is always covered, even when its dimensions are not an exact multiple
of the stride.

\paragraph{Sparse-mask handling.}
Fiber pixels are a small minority of total image area, so a large majority
of naively generated tiles are entirely background. During training (train
split only), tiles below a minimum foreground fraction (default 0.1\%) are
randomly kept with a fixed low probability (default 20\%) rather than
discarded outright -- this reduces the volume of uninformative background
tiles without eliminating the model's exposure to hard negatives. An
alternative dynamic mode exists (per-epoch resampling weighted toward the
model's current hardest negatives, i.e.\ background tiles it currently
misclassifies) but is not used in the production recipe.

\paragraph{Channel handling and standardization.}
The single grayscale channel is replicated to 3 channels so it is compatible
with ImageNet/MicroNet-pretrained encoders. After augmentation, tiles are
standardized with the ImageNet per-channel mean/std (matching the
distribution the pretrained encoder was fine-tuned on); this standardization
is applied identically in the training dataset and in the tiled-inference
code path.

\section{Data augmentation}

Applied to the training split only (validation/test are unaugmented), via
Albumentations:

\begin{table}[h]
\centering
\begin{tabular}{@{}llc@{}}
\toprule
Transform & Parameters & $p$ \\
\midrule
Horizontal flip & --- & 0.50 \\
Vertical flip & --- & 0.50 \\
$90^\circ$ rotation & --- & 0.50 \\
Affine & scale 0.75--1.35, translate $\pm5\%$, rotate $\pm15^\circ$ & 0.50 \\
Brightness/contrast jitter & $\pm0.15$ & 0.40 \\
Gaussian noise & $\sigma$ 0.02--0.08 (on the $[0,1]$ scale) & 0.30 \\
Gaussian blur & kernel 3--5 & 0.20 \\
\bottomrule
\end{tabular}
\caption{Training-time augmentation pipeline.}
\end{table}

The wide affine scale range specifically accounts for fiber pixel-width
varying with SEM magnification across the source dataset. Output is
re-clipped to $[0, 1]$ after augmentation, since brightness/noise transforms
can push values slightly outside the normalized range.

\section{Model architecture}

The segmentation backbone is a U-Net (\texttt{segmentation\_models\_pytorch})
with a ResNet-50 encoder. Encoder weights are initialized from NASA's
\textbf{MicroNet} pretrained-microscopy-models corpus -- a domain-specific
pretraining corpus of microscopy imagery, rather than natural-image ImageNet
-- under the hypothesis that microscopy-domain pretraining transfers better
to SEM textures than natural-image pretraining. The training/inference code
is architecture-agnostic: any \texttt{segmentation\_models\_pytorch} decoder
family (U-Net, U-Net++, FPN, DeepLabV3+) and encoder is swappable via
configuration without code changes, and standard ImageNet-pretrained or
from-scratch encoders are supported as alternatives to MicroNet.

The model outputs a single-channel raw logit map. Any built-in output
activation (sigmoid/softmax) present in third-party model definitions is
explicitly stripped, so sigmoid is applied only inside the loss, metric, and
inference code -- this ``raw logits'' contract is enforced uniformly
regardless of which backend produced the model.

\section{Loss function}

The loss is a weighted sum of three terms, computed on sigmoid probabilities
of the raw logits.

\paragraph{(a) Binary cross-entropy.} Standard pixelwise BCE.

\paragraph{(b) Focal Tversky loss.}
The soft Tversky index is

\begin{equation}
\mathrm{TI} = \frac{\mathrm{TP} + \varepsilon}{\mathrm{TP} + \alpha\cdot\mathrm{FP} + \beta\cdot\mathrm{FN} + \varepsilon}
\label{eq:tversky}
\end{equation}

with TP/FP/FN computed as probability-weighted sums (not binarized),
$\alpha = 0.3$, $\beta = 0.7$, $\varepsilon = 10^{-7}$. Weighting
$\beta > \alpha$ penalizes missed fiber pixels (false negatives) roughly
$2.3\times$ more heavily than false positives -- appropriate given that
under-segmenting a thin fiber is a more consequential error than mild
over-segmentation for downstream fiber counting. The focal form,

\begin{equation}
L_{FT} = (1 - \mathrm{TI})^{\gamma}, \quad \gamma = 0.75
\label{eq:focal-tversky}
\end{equation}

additionally up-weights harder (lower-TI) examples during training.

\paragraph{(c) Soft centerline-Dice (clDice) topology term.}
Weighted 0.3 in the total loss. A differentiable morphological skeleton is
computed for both the predicted probability map and the ground-truth mask
via iterative erosion (implemented as negated max-pooling, 5 iterations) and
re-dilation; clDice is the harmonic mean of (i) the fraction of the
predicted skeleton that lies inside the true mask and (ii) the fraction of
the true skeleton that lies inside the predicted mask. This term directly
rewards connected, unbroken fiber topology, addressing a failure mode of pure
overlap losses (BCE, Dice, Tversky), which tolerate a single fiber being
predicted as several disconnected fragments as long as total pixel overlap is
high.

Total training loss:

\begin{equation}
L = \mathrm{BCE} + L_{FT} + 0.3\cdot(1 - \mathrm{clDice})
\label{eq:total-loss}
\end{equation}

\section{Training procedure}

\begin{itemize}[nosep]
  \item \textbf{Optimizer:} AdamW, weight decay $10^{-4}$.
  \item \textbf{Differential learning rate:} base learning rate
        $2\times10^{-4}$ applied to the randomly-initialized decoder/head;
        the pretrained encoder receives this rate scaled by 0.1
        ($2\times10^{-5}$), so the transferred features are fine-tuned
        gently while the fresh decoder learns quickly.
  \item \textbf{LR schedule:} \texttt{ReduceLROnPlateau} (factor 0.5,
        patience 4 epochs) driven by the monitor metric
        (Section~\ref{sec:monitor}); cosine annealing is supported as an
        alternative.
  \item \textbf{Early stopping:} patience 20 epochs, minimum delta
        $10^{-3}$, on the same monitor metric.
  \item \textbf{Budget:} up to 400 epochs (early stopping typically halts
        training well before this ceiling).
  \item \textbf{Batch size:} 8 tiles of $512\times512$~px.
\end{itemize}

Stochastic Weight Averaging is supported but disabled in the production
recipe.

\section{Training-control metric}
\label{sec:monitor}

Checkpoint selection, early stopping, and LR scheduling are driven by a
\textbf{threshold-free} validation metric, \texttt{val/soft\_tversky} -- the
same Tversky formula as Equation~\ref{eq:tversky}, but accumulated as
probability-weighted TP/FP/FN sums over the \emph{entire} validation epoch
(not averaged per-batch, which would let empty tiles trivially score near
zero and deflate the true whole-image score on such sparse masks) and never
binarized at a threshold. This is a deliberate design choice: it decouples
training-time model selection from \texttt{train.threshold}, a deployment
decision that is not yet calibrated during training (see
Section~\ref{sec:calibration}). Hard, thresholded analogues of every metric
(\texttt{val/dice}, \texttt{val/tversky}, etc., at the eventual deployment
threshold) are logged in parallel for human-readable, deployment-realistic
reporting, but do not drive any training decision.

\section{Post-hoc threshold calibration}
\label{sec:calibration}

After training, the sigmoid-probability decision threshold is calibrated --
rather than left at an arbitrary default of 0.5 -- by sweeping it across
$\sim$99 values evaluated \textbf{only on the validation split} (never test,
to avoid leakage into the held-out evaluation set) and selecting the value
that maximizes a chosen metric (Dice by default). A precision--recall curve
and its AUC are reported for this sweep; precision--recall (rather than ROC)
is used because fiber masks are strongly background-dominated, and ROC-AUC
would look misleadingly high under that class imbalance -- the same
reasoning that motivates using Dice/IoU/Tversky over raw pixel accuracy
throughout this pipeline. The single resulting threshold is then applied
uniformly across every split at inference time.

\section{Inference}

Full-resolution images are tiled using the same patch size/stride as
training, so the effective receptive field and context seen by the model
matches training exactly. Overlapping tile probabilities are blended with a
2D Gaussian window (following the nnU-Net sliding-window scheme), which
down-weights each tile's less-reliable border pixels relative to its center
and smooths seams between adjacent tiles. Edge tiles are reflect-padded
(rather than zero-padded) so they are not artificially biased toward
``background'' by a black border.

Optional $8\times$ dihedral test-time augmentation (TTA) -- every combination
of flip/$90^\circ$ rotation, each inverse-transformed and
probability-averaged -- is available as a no-retrain accuracy gain
(empirically on the order of $+1$--$2$ Dice) at $8\times$ inference cost; it
is left off during iterative development and enabled only for final reported
predictions. The final binary mask is the blended probability map
thresholded at the calibrated value from Section~\ref{sec:calibration}.

\section{Post-processing}

Every predicted mask additionally passes through a fixed, classical
(non-learned) morphological cleanup, applied in this order:

\begin{enumerate}[nosep]
  \item Binary closing (disk radius 1) -- run \emph{first} so it can bridge
        small, genuine gaps in a single fiber before the next step would
        otherwise mistake the resulting fragments for separate noise specks.
  \item Binary opening (disabled by default, radius 0) -- available to
        smooth jagged edges, but disabled because it risks eroding away
        genuinely thin fibers.
  \item Removal of connected components $\leq 64$~px.
  \item Filling of interior holes $\leq 64$~px.
\end{enumerate}

Both the raw and post-processed masks are scored independently against
ground truth, and the per-image delta (post $-$ raw) for every metric is
reported alongside the absolute values, so the contribution of this
deterministic cleanup stage is always visible rather than assumed.

\section{Evaluation metrics}

For every image, the full confusion matrix (TP/FP/FN/TN, pixelwise) is
computed against ground truth, from which accuracy, precision, recall,
specificity, Dice, IoU, Tversky ($\alpha=0.3$, $\beta=0.7$ -- the same
asymmetric weighting used in the loss, so optimization and reporting stay
consistent), and F2 (a recall-weighted F-beta score) are derived. Metrics are
computed identically for raw and post-processed predictions, across every
split, and written to a single CSV for auditability.

\section{Experiment tracking and reproducibility}

Every training run logs its fully resolved configuration, per-step and final
metrics, and periodic/best qualitative prediction images to a local MLflow
tracking server, keyed by experiment name. The entire pipeline -- data
split, preprocessing, model, loss, training schedule, inference, and
post-processing -- is defined by one declarative YAML configuration per
experiment (Table~\ref{tab:config}), so reproducing a reported result
requires no code changes, only re-running training against the same config.

\section{Application context}

The source SEM images originate from air-filter samples collected for
occupational asbestos-exposure monitoring. Segmented fiber masks feed
downstream quantification -- filter loading, WHO-criteria fiber counting,
and airborne fiber concentration / time-weighted-average exposure
reporting -- that otherwise relies entirely on manual microscopist counting
under magnification. The model's intended role is to \textbf{assist and
accelerate}, not silently replace, that expert workflow: the held-out test
split has no influence on threshold calibration
(Section~\ref{sec:calibration}), and companion tooling in the same
repository supports manual per-image quality review (mask quality, filter
loading level, and asbestos presence) so segmentation output remains
auditable against expert judgment before being used in any exposure-limit-
relevant reporting.

\begin{table}[h]
\centering
\begin{tabular}{@{}ll@{}}
\toprule
Setting & Value \\
\midrule
Architecture / encoder & U-Net, ResNet-50, MicroNet-pretrained \\
Input channels & 3 (grayscale replicated) \\
Patch size / stride & 512 / 384 px (25\% overlap) \\
Batch size & 8 \\
Split (train/val/test), seed & 70/15/15, seed 42 \\
Loss & BCE + Focal Tversky ($\alpha$=0.3, $\beta$=0.7, $\gamma$=0.75) + 0.3$\times$clDice (5 iters) \\
Optimizer & AdamW, lr $2\times10^{-4}$ (decoder) / $2\times10^{-5}$ (encoder), wd $10^{-4}$ \\
LR schedule & ReduceLROnPlateau, factor 0.5, patience 4 \\
Early stopping & patience 20, min $\Delta$ $10^{-3}$ \\
Max epochs & 400 \\
Monitor metric & \texttt{val/soft\_tversky} (threshold-free, maximize) \\
Inference blending & Gaussian window, reflect padding \\
Test-time augmentation & $8\times$ dihedral (final predictions only) \\
Post-processing & close (r=1) $\rightarrow$ open (r=0, off) $\rightarrow$ remove $\leq$64px $\rightarrow$ fill $\leq$64px \\
\bottomrule
\end{tabular}
\caption{Production configuration.}
\label{tab:config}
\end{table}

\end{document}
